\subsection{\emph{Prompts} de Pequena Dimensão}

\paragraph{S0} What is the capital of France?

\paragraph{S1} Which country is the largest producer of cork?

\paragraph{S2} What is the chemical symbol of gold?

\paragraph{S3} Who painted "The Scream"?

\paragraph{S4} What is the second planet counting from the Sun?

\subsection{\emph{Prompts} de Média Dimensão}

\paragraph{M0} Explain the key differences between machine learning and deep learning, and provide
               an example of each.

\paragraph{M1} Explain the difference between speed and velocity, and give an example to illustrate
               each concept.

\paragraph{M2} Explain the concept of supply and demand, and how it affects the price of goods.

\paragraph{M3} Explain the difference between matrix-matrix multiplication and matrix-vector
               multiplication. Give a practical application for each one.

\paragraph{M4} Explain the difference in chemical composition between acid and alkaline solutions.
               Give examples of both.

\subsection{\emph{Prompts} de Grande Dimensão}

\paragraph{L0} The transformer architecture, introduced in the paper ``Attention is All You Need''
               by Vaswani et al. in 2017, has become the foundation for most modern large language
               models. Unlike recurrent neural networks, which process sequences sequentially,
               transformers process all tokens in parallel using self-attention mechanisms. This
               allows better parallelisation during training and has proven more effective at
               capturing long-range dependencies in text. When deploying transformers for inference
               on CPU, the main bottlenecks are the model’s large memory footprint and the repeated
               loading of weights during the decoding phase. Based on this context, what are the
               main challenges when running large transformer models on CPU-only systems, and how
               might quantisation help address these challenges?

\paragraph{L1} The Paris Agreement aims to limit global warming to well below 2°C above
               pre-industrial levels, with efforts to cap it at 1.5°C. Achieving this requires
               rapid decarbonization, particularly in the energy sector. While solar and wind power
               have seen dramatic cost reductions, their intermittency remains a major challenge for
               grid stability. Energy storage solutions, such as lithium-ion batteries, are
               improving but still face limitations in scalability, cost, and environmental impact.
               Based on this context, what are the main obstacles to achieving a fully renewable
               energy grid, and how might advances in green hydrogen or next-generation battery
               technologies help address these challenges?

\paragraph{L2} Dark matter, which constitutes approximately 27\% of the universe's mass-energy
               content, does not emit, absorb, or reflect light, making it invisible to traditional
               telescopes. Its existence is inferred from gravitational effects on visible matter,
               such as the rotation curves of galaxies and the behavior of galaxy clusters. Despite
               decades of research, the particle nature of dark matter remains unknown, with leading
               candidates including weakly interacting massive particles (WIMPs) and axions. Direct
               detection experiments, such as those using underground detectors or space-based
               observatories, have yet to yield conclusive evidence. Based on this context, what are
               the main challenges in detecting and identifying dark matter, and how might
               next-generation experiments or alternative theoretical models help resolve these
               challenges?

\paragraph{L3} Antimicrobial resistance (AMR) is a growing global health crisis, with bacteria,
               viruses, and fungi evolving to resist existing drugs. Overuse and misuse of
               antibiotics in healthcare and agriculture have accelerated this process, leading to
               infections that are increasingly difficult to treat. Low- and middle-income countries
               are particularly vulnerable due to limited access to diagnostics and alternative
               treatments. In this context, what are the main challenges in combating antimicrobial
               resistance?

\paragraph{L4} Portugal's higher education system is bifurcated into universities, which focus on
               theoretical and research-oriented education, and polytechnics, which emphasize
               applied, vocational training. While polytechnics play a crucial role in addressing
               labor market needs, they often face challenges such as lower prestige, limited
               research funding, and fewer opportunities for academic progression. With the growing
               demand for skilled technicians and the push for lifelong learning, how can
               polytechnics in Portugal enhance their status and better align their offerings with
               evolving industry requirements?
